\section*{S1 Software and settings of the competitors}
\label{supp:software}

TCLUST-REG is the function \texttt{tclustreg} of the FSDA toolbox, a MATLAB toolbox. We ran it
unmodified (GitHub commit \texttt{37f998e} of 24 September 2026) under GNU Octave~9.4 with the
\texttt{statistics} package ($1.6.7$ for the simulations, $1.6.3$ for the data of Section~5),
with three one-line substitutes for MATLAB functions that Octave lacks: a version check that
returns false, \texttt{coder.target}, which returns true for \texttt{"MATLAB"}, and
\texttt{convertStringsToChars}, which returns its arguments; they are in the repository. The runs
with adaptive second-level trimming (Section~S6) needed a fourth, for
\texttt{any(x,'all')}, which Octave's \texttt{any} does not accept.

TCLUST-REG was then run again in MATLAB (R2026a Update~5 in MATLAB Online, FSDA at the same
commit, no substitutes) on $80$ data sets, ten replications of each of the eight design cells of
Table~\ref{tab:matlab}, each at both levels with $1000$ starts and $50$ refinement steps. MATLAB's
random number generator differs from Octave's, so the two programs drew different starts. They
returned the same coefficients, to the ten significant digits both programs write and up to the
labelling of the groups, in $68$ of the $160$ fits; in the other $92$ the largest coefficient
difference exceeded $2\times10^{-3}$ of the largest coefficient. The Octave runs did not record the objective, so each fit was also scored by one criterion computed from its
coefficients: concentration steps at the fit's level with the ratio of the group variances bounded
by $12$, then the trimmed classification log-likelihood. On the MATLAB fits this score equalled
the objective that MATLAB reports to within $0.5$ in $137$ of the $160$ fits and differed from it
by up to $13$ in the rest. MATLAB's fit scored higher by more than
$0.5$ in $28$ fits and Octave's in $30$; the rest were within $0.5$. With the margin raised to $2$ the counts were $21$ and $22$, and
with $5$ they were $17$ and $12$. Neither program found better
optima more often, and the differences are those of two random searches. After reweighting and
the recovery step the mean accuracy over the $160$ fits was $0.836$ in MATLAB and $0.830$ in
Octave. Cell means differed by up to $0.09$, and by more than $0.05$ only in cells where no fit agreed. MATLAB Online took a median of $1.2$ seconds per fit and Octave, on our cluster, $41$ seconds.
The two ran on different machines, so the ratio of about $36$ mixes interpreter and hardware; the
times of TCLUST-REG in Table~\ref{tab:time}, measured on the cluster, overstate its cost in MATLAB
by a factor of that order.

The trimmed cluster-weighted model, TLE, CTLE, the bisquare mixture and the Laplace mixture come
from the R package \texttt{RobMixReg}~1.1.3; the trimmed cluster-weighted model was given $50$
random starts and $20$ concentration steps, more than its defaults, and the Laplace mixture is
\texttt{mixLp} with \texttt{nit}$=20$. The bisquare mixture \texttt{mixlinrb\_bi} raised an error
in our R session, because its single-start routine fits \texttt{lm(formula, data, weights = w)}
with a formula whose environment does not contain \texttt{w}; we ran a copy of the two package
functions in which only this line sets the formula environment first (\texttt{bisq\_patched} in
\texttt{r\_methods.R}). The computation is otherwise the package's. No package offers the contaminated normal and the
noise-component mixtures of regressions, so the implementations are ours. The contaminated
normal mixture is fitted by an ECM algorithm from $50$ starts, each from $K$ random elemental
lines with the units far from all of them starting as bad points, with the inflation started at
$20$ and the ratio of the group variances bounded by $12$; starts from random partitions, which
we tried first, were captured by the outliers. The noise-component mixture is fitted by EM from
$10$ random partitions, the best continued after $10$ iterations, with the ratio of the component
variances bounded by $12$, as in TCLUST-REG. Table~\ref{tab:S21} repeats the three mixtures with
more starts.

A simple trimmed alternation of our own at the same five levels, with the reweighting step,
completes the set that the comparison of Plan~1 used; it is a home-made
baseline, so it appears in the Supplement only, although it is the best competitor in some
cells of Table~\ref{tab:S1}.
